% Gerado por ferramentas/gerar.py; não edite à mão.
\section*{Exercícios}\markright{Exercícios}
\addcontentsline{toc}{section}{Exercícios}
Os pontos antes de cada número indicam a dificuldade: \nivel{1} fácil, \nivel{2} médio, \nivel{3} difícil e \nivel{4} desafio. Os exercícios marcados com \estrela{} têm solução completa, com animação do gráfico, no repositório do livro: \url{https://github.com/AleffAzevedo/calculo-2-pt/blob/main/exercicios/cap1/README.md}. As respostas estão no fim do capítulo.

\begin{itemize}[leftmargin=5.2em,labelsep=0.6em,itemsep=4pt]
\item[\exnum{1}{1.1}] Dê a ordem de cada equação e diga se ela é linear: (a) $y''+3y'-y=\sen x$; (b) $(y')^2+y=x$; (c) $y'''=y\,y'$.
\item[\exnum{1}{1.2}] Verifique que $y=C_1e^{2x}+C_2e^{-x}$ é solução de $y''-y'-2y=0$ para quaisquer constantes $C_1$, $C_2$.
\item[\exnum{1}{1.3}] Para quais valores de $r$ a função $y=e^{rx}$ é solução de $y''+y'-6y=0$?
\item[\exnum{1}{1.4}] Resolva o PVI $y'=x\,e^{x^2}$, $y(0)=1$.
\item[\exnum{1}{1.5}] Resolva o PVI $y''=6x$, $y(0)=1$, $y'(0)=2$.
\item[\exnum{1}{1.6}] Uma pedra é lançada para cima, do chão, a $20$ m/s ($g=9{,}8$ m/s$^2$). Escreva $v(t)$ e $s(t)$ e encontre a altura máxima.
\item[\exnum{1}{1.7}] Encontre e classifique os equilíbrios de $y'=y^2-4$.
\item[\exnum{1}{1.8}] Encontre e classifique os equilíbrios de $y'=y(y-3)^2$.
\item[\exnum{1}{1.9}] Faça dois passos do método de Euler com $h=0{,}5$ para $y'=2x-y$, $y(0)=1$, e compare com a solução exata $y=2x-2+3e^{-x}$ em $x=1$.
\item[\exnum{1}{1.10}] Resolva $y'=\dfrac yx$ ($x>0$) por separação de variáveis.
\item[\exnum{1}{1.11}] Resolva $y'=e^{x-y}$.
\item[\exnum{1}{1.12}] Resolva $y'=1+y^2$, $y(0)=0$. Em que intervalo a solução existe?
\item[\exnum{1}{1.13}] Resolva $y'=x\,y^2$, $y(0)=1$, e diga onde a solução deixa de existir.
\item[\exnum{1}{1.14}] Um corpo a $80$\textdegree C é deixado num ambiente a $20$\textdegree C, com $T'=-0{,}1(T-20)$ ($t$ em min). Encontre $T(t)$ e quando ele chega a $40$\textdegree C.
\item[\exnum{1}{1.15}] Para $P'=0{,}5P\big(1-\tfrac P{200}\big)$, dê os equilíbrios, a capacidade de suporte e o valor de $P$ em que o crescimento é mais rápido (com a taxa nesse ponto).
\item[\exnum{1}{1.16}] Com os dados do exercício anterior e $P(0)=20$, calcule $P(4)$.
\item[\exnum{1}{1.17}] Resolva $y'+2y=4$ e descreva o comportamento das soluções quando $x\to\infty$.
\item[\exnum{1}{1.18}] Resolva $y'+\dfrac yx=3x$ ($x>0$) pelo fator integrante.
\item[\exnum{1}{1.19}] Resolva o PVI $y'-3y=e^{2x}$, $y(0)=0$.
\item[\exnum{1}{1.20}] Um tanque tem $100$ L de água pura. Entra salmoura com $0{,}2$ kg/L a $5$ L/min e sai mistura a $5$ L/min. Encontre o sal $u(t)$ e quanto há após $20$ min.
\item[\exnum{2}{1.21}] \estrela\ \textbf{Um campo que espalha as soluções.} Para $y'=y-x$, esboce o campo de direções, encontre a solução geral e mostre que só uma
solução não se afasta da reta $y=x+1$. Compare $y(0)=1$ e $y(0)=1{,}1$ em $x=5$.
\item[\exnum{2}{1.22}] \estrela\ \textbf{Euler e o número $e$.} Aplique o método de Euler a $y'=y$, $y(0)=1$, em $[0,1]$ com $n$ passos de tamanho $h=\frac1n$.
Mostre que a aproximação de $y(1)$ é $\big(1+\frac1n\big)^n$ e calcule-a para $n=1,2,4,10,100$.
\item[\exnum{2}{1.23}] \estrela\ \textbf{Soluções que acabam.} Resolva $y'=-\dfrac xy$, $y(0)=2$. Em que intervalo a solução existe? O que acontece nas pontas?
\item[\exnum{2}{1.24}] \estrela\ \textbf{A hora da morte.} Um corpo é encontrado a $30$\textdegree C numa sala a $20$\textdegree C; uma hora depois está a
$28$\textdegree C. Supondo a lei de Newton e temperatura de $37$\textdegree C no momento da morte, há quanto tempo
ocorreu a morte quando o corpo foi encontrado?
\item[\exnum{2}{1.25}] \estrela\ \textbf{Resolvendo a logística.} Resolva $P'=rP\big(1-\frac PK\big)$, $P(0)=P_0$, por separação de variáveis e frações parciais.
Para $r=0{,}2311$, $K=4000$ e $P_0=900$, quando a população cresce mais rápido?
\item[\exnum{2}{1.26}] \estrela\ \textbf{Velocidade terminal.} Um paraquedista cai com $v'=9{,}8-0{,}2v$, $v(0)=0$ (SI). Encontre $v(t)$, a distância percorrida
$s(t)$ e quanto tempo leva para atingir $95\%$ da velocidade terminal.
\item[\exnum{2}{1.27}] Classifique os equilíbrios de $y'=(y-1)(y-4)$ e esboce algumas soluções.
\item[\exnum{2}{1.28}] Use Euler com $h=0{,}25$ (quatro passos) para estimar $y(1)$ em $y'=x^2-y$, $y(0)=1$. A solução exata é $y=x^2-2x+2-e^{-x}$.
\item[\exnum{2}{1.29}] Resolva $y'=3x^2(y+1)$, $y(0)=0$.
\item[\exnum{2}{1.30}] Resolva $y'=\dfrac{\cos x}{y}$, $y(0)=2$.
\item[\exnum{2}{1.31}] Resolva $\dfrac{dy}{dx}=\dfrac{1+y^2}{1+x^2}$, $y(0)=1$, e simplifique a resposta.
\item[\exnum{2}{1.32}] Resolva $y'=y\ln y$, $y(0)=e$.
\item[\exnum{2}{1.33}] Resolva $y'+2xy=x$, $y(0)=0$.
\item[\exnum{2}{1.34}] Resolva $y'+y\tg x=\sec x$, $y(0)=1$, em $|x|<\tfrac\pi2$.
\item[\exnum{2}{1.35}] Resolva $xy'-2y=x^3e^x$ ($x>0$).
\item[\exnum{2}{1.36}] Escreva $P'=0{,}02P-0{,}00004P^2$ na forma logística e identifique $r$ e $K$.
\item[\exnum{2}{1.37}] Uma população logística tem $K=1000$, $P(0)=100$ e atinge $500$ em $t=10$. Encontre $r$.
\item[\exnum{2}{1.38}] Estime $y(1)$ para $y'=y$, $y(0)=1$, com Euler de passo $h=0{,}5$ e $h=0{,}25$. Compare os erros.
\item[\exnum{2}{1.39}] Um bolo sai do forno a $180$\textdegree C para uma sala a $25$\textdegree C e está a $100$\textdegree C após $10$ min. Quando estará a $50$\textdegree C?
\item[\exnum{2}{1.40}] Verifique que $y=\dfrac1{1+Ce^{-x}}$ resolve $y'=y(1-y)$ e encontre a solução com $y(0)=\tfrac14$.
\item[\exnum{3}{1.41}] \estrela\ \textbf{Logística com colheita.} Uma população de peixes segue $P'=0{,}4P\big(1-\frac P{1000}\big)-60$ (pesca constante de 60 por ano).
Encontre e classifique os equilíbrios. O que acontece se $P(0)=150$? Qual a maior pesca sustentável?
\item[\exnum{3}{1.42}] \estrela\ \textbf{Mistura com volume variável.} Um tanque tem $100$ L de água pura. Entra salmoura com $0{,}5$ kg/L a $3$ L/min e sai a mistura a
$2$ L/min. Encontre o sal $u(t)$ e a concentração quando o volume dobrar.
\item[\exnum{3}{1.43}] \estrela\ \textbf{Uma equação de Bernoulli.} Resolva $y'+y=y^2$ com a substituição $v=\frac1y$. Descreva as soluções com $y(0)=\tfrac13$, $y(0)=1$ e $y(0)=2$.
\item[\exnum{3}{1.44}] \estrela\ \textbf{Trajetórias ortogonais.} Encontre as curvas que cortam todas as parábolas $y=Cx^2$ em ângulo reto.
\item[\exnum{3}{1.45}] \estrela\ \textbf{Esvaziando um tanque.} Pela lei de Torricelli, o nível de um tanque cilíndrico que esvazia por um furo no fundo satisfaz
$h'=-0{,}1\sqrt h$ ($h$ em m, $t$ em min). Com $h(0)=4$, encontre $h(t)$, o tempo para esvaziar e o tempo
para o nível cair à metade.
\item[\exnum{3}{1.46}] Mostre que o PVI $y'=3y^{2/3}$, $y(0)=0$ tem pelo menos duas soluções. Por que isso não contradiz a intuição do campo de direções?
\item[\exnum{3}{1.47}] Resolva $y'=y^2-1$, $y(0)=0$, usando frações parciais.
\item[\exnum{3}{1.48}] Resolva $y'=\dfrac{x^2}{1-y^2}$ (solução implícita).
\item[\exnum{3}{1.49}] Resolva $y'+y=\sen x$.
\item[\exnum{3}{1.50}] Resolva $y'-\dfrac yx=x\sen x$ ($x>0$).
\item[\exnum{3}{1.51}] Resolva $(x+y^2)\,dy=y\,dx$ tratando $x$ como função de $y$.
\item[\exnum{3}{1.52}] Aplique Euler a $y'=-2y$, $y(0)=1$, com $h=1{,}5$. O que acontece? Para quais $h$ as aproximações decaem como a solução exata?
\item[\exnum{3}{1.53}] Mostre, derivando a equação logística, que $P''=rP'\big(1-\tfrac{2P}K\big)$ e conclua que a inflexão ocorre em $P=\tfrac K2$.
\item[\exnum{3}{1.54}] O modelo de Gompertz é $P'=rP\ln\frac KP$. Resolva-o com a substituição $u=\ln P$.
\item[\exnum{3}{1.55}] Um medicamento é infundido a $2$ mg/h e eliminado a uma taxa de $10\%$ da quantidade por hora. Com $A(0)=0$, encontre $A(t)$ e quando se atinge $90\%$ do nível de equilíbrio.
\item[\exnum{3}{1.56}] Um investimento rende $5\%$ ao ano (capitalização contínua) e recebe depósitos contínuos de R\$ $1200$ por ano. Com $S(0)=0$, quanto há após $20$ anos?
\item[\exnum{3}{1.57}] Com resistência quadrática, $v'=9{,}8-0{,}098v^2$, $v(0)=0$. Encontre $v(t)$ e a velocidade terminal.
\item[\exnum{3}{1.58}] Resolva $y'=(x+y)^2$ com a substituição $u=x+y$.
\item[\exnum{3}{1.59}] Resolva a equação homogênea $y'=\dfrac{x^2+y^2}{xy}$ ($x>0$) com $v=\frac yx$.
\item[\exnum{3}{1.60}] Uma doença se espalha numa cidade de $5000$ pessoas segundo $y'=0{,}0002\,y(5000-y)$, com $y(0)=10$. Quando metade estará infectada?
\item[\exnum{4}{1.61}] Mostre que toda solução de $y'+p(x)y=q(x)$ (com $p$, $q$ contínuas) é $y=e^{-\int p}\Big(\int e^{\int p}\,q\,dx+C\Big)$ e que o PVI com $y(x_0)=y_0$ tem exatamente uma solução.
\item[\exnum{4}{1.62}] Encontre as trajetórias ortogonais da família de circunferências $x^2+y^2=2cx$.
\item[\exnum{4}{1.63}] (Clepsidra) Pela lei de Torricelli, $A(h)\,h'=-k\sqrt h$. Que forma deve ter um recipiente de revolução para que o nível desça a taxa constante?
\item[\exnum{4}{1.64}] Um cabo suspenso satisfaz $y''=\tfrac1a\sqrt{1+(y')^2}$. Com $p=y'$, resolva e mostre que $y=a\cosh\frac xa$ (com $y'(0)=0$, $y(0)=a$).
\item[\exnum{4}{1.65}] (Riccati) Verifique que $y_1=\tfrac1x$ resolve $y'=y^2-\tfrac2{x^2}$ e encontre a solução geral com $y=y_1+\tfrac1v$.
\end{itemize}
\subsection*{Respostas}\markright{Respostas}
{\small\begin{multicols}{2}\raggedright
\begin{itemize}[leftmargin=2.7em,labelsep=0.5em,itemsep=1pt]
\item[\textbf{1.1}] (a) ordem 2, linear; (b) ordem 1, não linear; (c) ordem 3, não linear.
\item[\textbf{1.2}] $4C_1e^{2x}+C_2e^{-x}-(2C_1e^{2x}-C_2e^{-x})-2(C_1e^{2x}+C_2e^{-x})=0$.
\item[\textbf{1.3}] $r^2+r-6=0$: $r=2$ ou $r=-3$.
\item[\textbf{1.4}] $y=\tfrac12e^{x^2}+\tfrac12$.
\item[\textbf{1.5}] $y=x^3+2x+1$.
\item[\textbf{1.6}] $v=20-9{,}8t$, $s=20t-4{,}9t^2$; máximo em $t\approx2{,}04$ s, $s\approx20{,}4$ m.
\item[\textbf{1.7}] $y=-2$ estável; $y=2$ instável.
\item[\textbf{1.8}] $y=0$ instável; $y=3$ semiestável.
\item[\textbf{1.9}] $y_1=0{,}5$, $y_2=0{,}75$; exato $3/e\approx1{,}104$.
\item[\textbf{1.10}] $y=Cx$.
\item[\textbf{1.11}] $e^y=e^x+C$, ou $y=\ln(e^x+C)$.
\item[\textbf{1.12}] $y=\tg x$, para $-\tfrac\pi2<x<\tfrac\pi2$.
\item[\textbf{1.13}] $y=\dfrac2{2-x^2}$; explode em $x=\pm\sqrt2$.
\item[\textbf{1.14}] $T=20+60e^{-0{,}1t}$; $t=10\ln3\approx11{,}0$ min.
\item[\textbf{1.15}] $P=0$ e $P=200=K$; $P=100$, com $P'=25$.
\item[\textbf{1.16}] $P(4)=\dfrac{4000e^2}{180+20e^2}\approx90{,}2$.
\item[\textbf{1.17}] $y=2+Ce^{-2x}\to2$.
\item[\textbf{1.18}] $\mu=x$; $y=x^2+\dfrac Cx$.
\item[\textbf{1.19}] $y=e^{3x}-e^{2x}$.
\item[\textbf{1.20}] $u=20(1-e^{-t/20})$; $u(20)=20(1-e^{-1})\approx12{,}6$ kg.
\item[\textbf{1.21}] $y=x+1+Ce^x$; só $C=0$; $y(5)=6$ contra $y(5)\approx20{,}8$.
\item[\textbf{1.22}] $2;\ 2{,}25;\ 2{,}441;\ 2{,}594;\ 2{,}705$, tendendo a $e\approx2{,}718$.
\item[\textbf{1.23}] $y=\sqrt{4-x^2}$, para $-2<x<2$; a inclinação tende a $\mp\infty$ nas pontas.
\item[\textbf{1.24}] Cerca de $2{,}38$ h, isto é, $2$ h $23$ min.
\item[\textbf{1.25}] $P=\dfrac{KP_0e^{rt}}{K-P_0+P_0e^{rt}}$; em $t=\frac1r\ln\frac{K-P_0}{P_0}\approx5{,}35$ anos, com $P=2000$.
\item[\textbf{1.26}] $v=49(1-e^{-0{,}2t})$; $s=49t-245(1-e^{-0{,}2t})$; $t=5\ln20\approx15$ s.
\item[\textbf{1.27}] $y=1$ estável; $y=4$ instável.
\item[\textbf{1.28}] $0{,}75;\ 0{,}578;\ 0{,}496;\ 0{,}513$; exato $1-e^{-1}\approx0{,}632$.
\item[\textbf{1.29}] $y=e^{x^3}-1$.
\item[\textbf{1.30}] $y=\sqrt{2\sen x+4}$.
\item[\textbf{1.31}] $\arctg y=\arctg x+\tfrac\pi4$, ou $y=\dfrac{1+x}{1-x}$.
\item[\textbf{1.32}] $y=e^{e^x}$.
\item[\textbf{1.33}] $y=\tfrac12\big(1-e^{-x^2}\big)$.
\item[\textbf{1.34}] $\mu=\sec x$; $y=\sen x+\cos x$.
\item[\textbf{1.35}] $y=x^2(e^x+C)$.
\item[\textbf{1.36}] $P'=0{,}02P(1-P/500)$: $r=0{,}02$, $K=500$.
\item[\textbf{1.37}] $r=\tfrac{\ln9}{10}\approx0{,}220$.
\item[\textbf{1.38}] $2{,}25$ e $2{,}441$; erros $0{,}47$ e $0{,}28$ (cerca da metade).
\item[\textbf{1.39}] $t\approx25{,}1$ min.
\item[\textbf{1.40}] $C=3$: $y=\dfrac1{1+3e^{-x}}$.
\item[\textbf{1.41}] $P\approx183{,}8$ (instável) e $P\approx816{,}2$ (estável); com $P(0)=150$ a população se extingue; pesca máxima $100$ por ano.
\item[\textbf{1.42}] $u=\tfrac12(100+t)-\dfrac{500\,000}{(100+t)^2}$; em $t=100$: $87{,}5$ kg, concentração $0{,}4375$ kg/L.
\item[\textbf{1.43}] $y=\dfrac1{1+Ce^x}$; tende a $0$, fica em $1$, e explode em $x=\ln2$, respectivamente.
\item[\textbf{1.44}] As elipses $x^2+2y^2=K$.
\item[\textbf{1.45}] $h=(2-0{,}05t)^2$; esvazia em $40$ min; metade do nível em $\approx11{,}7$ min.
\item[\textbf{1.46}] $y\equiv0$ e $y=x^3$; $f(y)=3y^{2/3}$ não tem derivada limitada perto de $y=0$.
\item[\textbf{1.47}] $y=\dfrac{1-e^{2x}}{1+e^{2x}}=-\tgh x$.
\item[\textbf{1.48}] $3y-y^3=x^3+C$.
\item[\textbf{1.49}] $y=\tfrac12(\sen x-\cos x)+Ce^{-x}$.
\item[\textbf{1.50}] $y=x(C-\cos x)$.
\item[\textbf{1.51}] $\dfrac{dx}{dy}-\dfrac xy=y$: $x=y^2+Cy$.
\item[\textbf{1.52}] $y_{n+1}=(1-2h)y_n=(-2)^n$ oscila e diverge; é preciso $|1-2h|<1$, isto é, $0<h<1$.
\item[\textbf{1.53}] Derive $rP-\tfrac rKP^2$ pela regra da cadeia.
\item[\textbf{1.54}] $u'=r(\ln K-u)$: $P=K\exp\big(Ce^{-rt}\big)$.
\item[\textbf{1.55}] $A=20(1-e^{-0{,}1t})$; $t=10\ln10\approx23$ h.
\item[\textbf{1.56}] $S=24\,000(e^{0{,}05t}-1)$; $S(20)\approx$ R\$ $41\,239$.
\item[\textbf{1.57}] $v=10\tgh(0{,}98t)$; terminal $10$ m/s.
\item[\textbf{1.58}] $u'=1+u^2$: $y=\tg(x+C)-x$.
\item[\textbf{1.59}] $xv'=\tfrac1v$: $y^2=2x^2\ln x+Cx^2$.
\item[\textbf{1.60}] Logística com $r=1$, $K=5000$: $t=\ln499\approx6{,}2$ dias.
\item[\textbf{1.61}] Se $y$ é solução, $(\mu y)'=\mu q$ com $\mu=e^{\int p}>0$; integre e use $y(x_0)$ para fixar $C$.
\item[\textbf{1.62}] $y'=\tfrac{y^2-x^2}{2xy}$ para a família; as ortogonais são $x^2+y^2=2ky$.
\item[\textbf{1.63}] $A=\pi x^2$ proporcional a $\sqrt h$: $x\propto h^{1/4}$, ou $h=cx^4$.
\item[\textbf{1.64}] $\int\tfrac{dp}{\sqrt{1+p^2}}=\tfrac xa$: $p=\senh\tfrac xa$, logo $y=a\cosh\tfrac xa$.
\item[\textbf{1.65}] $v'+\tfrac2xv=-1$: $y=\dfrac1x+\dfrac{3x^2}{C-x^3}$.
\end{itemize}\end{multicols}}
